\section{Authenticating Native MLE Requests}

The ideal authentication interface \(\mathcal F_{\mathrm{MLEAuth}}\) was defined
in Figure~\ref{fig:func-mle-auth}.  This section explains the implementation
shape used later by Blaze.  The native RAA compiler decides which tensors exist
and which evaluations are needed.  The authentication backend should not know
that RAA structure.  It should only authenticate claims of the form
\[
  \MLE(h)(r)=v
\]
for committed multilinear tensors \(h\).

Blaze implements this by batching the field-point claims and then using an RMLE
protocol that realizes the appropriate language-proximity functionality, with
BaseFold as the concrete example~[Blaze24,BaseFold23].  The batching step
follows the HyperPlonk-style MLE evaluation batching used by
Blaze~[CBBZ23,Blaze24].

The batching step is itself just sumcheck.  Suppose the verifier has already
fixed one committed tensor \(A\in\F^{\B^M}\), and wants to authenticate \(t\)
claims, where \([t]=\{1,\ldots,t\}\):
\[
  \MLE(A)(x_i)=v_i,
  \qquad x_i\in\F^M,\quad v_i\in\F,\quad i\in[t].
\]
For every point \(x\in\F^M\), the MLE identity says
\[
  \MLE(A)(x)=\sum_{b\in\B^M}A[b]\chi_b(x).
\]
After the verifier samples a batching challenge \(\gamma\), the \(t\) claims
are folded into the single Boolean-sum claim
\[
  \sum_{b\in\B^M}\MLE(A)(b)W_\gamma(b)
  =
  \sum_{i=1}^t\gamma^{i-1}v_i,
  \qquad
  W_\gamma(X)=\sum_{i=1}^t\gamma^{i-1}\chi_X(x_i).
\]
Here \(W_\gamma:\F^M\to\F\) is public once the request list and \(\gamma\) are
known.  If all \(t\) requested claims are true, then this Boolean-sum claim is
true exactly.  If at least one requested claim is false, then the folded claim
can still be true only when
\[
  \sum_{i=1}^t \gamma^{i-1}\bigl(\MLE(A)(x_i)-v_i\bigr)=0.
\]
For a fixed committed tensor \(A\) and fixed false request list, this is a
nonzero polynomial in \(\gamma\) of degree at most \(t-1\).  That is the
random-linear-combination loss, bounded by \((t-1)/|\F|\) by
Schwartz--Zippel.

Conditioned on the folded Boolean-sum claim being false, ordinary sumcheck
soundness is the next barrier.  Sumcheck on the product
\(\MLE(A)(X)W_\gamma(X)\) reduces the Boolean sum to one endpoint equation
\[
  \MLE(A)(r_{\mathrm{batch}})W_\gamma(r_{\mathrm{batch}})
  =
  T_{\mathrm{batch}}.
\]
The prover then supplies \(v_{\mathrm{batch}}\), the claimed value of
\(\MLE(A)(r_{\mathrm{batch}})\), and the backend authenticates that one MLE
claim.
The endpoint factor \(W_\gamma(r_{\mathrm{batch}})\) deserves a small warning.
When it is nonzero, the terminal equation pins down the required backend value.
When it is zero, the terminal equation no longer determines
\(v_{\mathrm{batch}}\); the transcript is then handled by the preceding
random-combination and sumcheck soundness analysis rather than by the final
backend opening.  One can conservatively count this as an endpoint degeneracy
event in the batching reduction.  When \(W_\gamma\) is not the zero polynomial,
Schwartz--Zippel bounds this event by \(M/|\F|\) over
\(r_{\mathrm{batch}}\).  It is not a separate commitment assumption.
The notation \(\chi_X(x_i)\) means the same equality polynomial as
\(\chi_b(x_i)\), but with the Boolean address \(b\) replaced by the field
variable \(X\).  Thus \(W_\gamma\) is a public multilinear polynomial in \(X\).

The real protocol also needs coordinate commitments, because the backend cannot
store all tensors in an ideal box.  These coordinate handles still use an
\(\mathsf{Open}\) API for Boolean-indexed entries, while MLE field-point claims
use \(\mathsf{Eval}\).  A real commitment phase therefore receives two pieces
for each oracle:
\[
  h_i\in\F^{\B^{d_i}}
  \qquad\text{and}\qquad
  C_i,
\]
where \(C_i\) is a coordinate-opening handle for \(h_i\).  The handle may be an
ordinary Merkle message-oracle handle, a composite message-oracle handle, or a
virtual message-oracle handle.  The implementation below treats all handles
opaquely.

The implementation uses one single-tensor backend subroutine.  Its job is not
to know where the tensor came from.  It receives a tensor
\(A\in\F^{\B^M}\) together with a coordinate-opening handle \(C_A\) for
Boolean entries of \(A\), returns an opaque backend handle, and later
authenticates one claim \(\MLE(A)(r)=v\).  The Blaze-facing API is therefore
\[
  \BackendRMLE.\mathsf{Commit}(A,C_A)\to A_{\mathrm{auth}},
  \qquad
  \BackendRMLE.\mathsf{Eval}(A_{\mathrm{auth}},r,v)\to\{0,1\}.
\]
The backend may internally compute an encoded word
\(\pi_M=\Enc_M(A)\), commit only to its non-systematic part, and route
systematic coordinates through \(C_A\).  That internal codeword is not an
object that the Blaze compiler inspects; it exists to define the RMLE relation
and the proximity check.

The backend handle \(A_{\mathrm{auth}}\) is therefore the boundary.  It stores
whatever commitment state the backend needs: the systematic coordinate handle
\(C_A\), any backend-owned parity commitments, routing data, dimensions, and
the public code parameters.  The evaluation phase uses that state to run the
backend's ordinary RMLE evaluation proof.  We do not expose a separate
``precommitted BaseFold'' API to Blaze.

This proximity check is doing real work.  A few authenticated coordinate
openings do not, by themselves, prove that a received word is a valid backend
encoding.  The foldable-code distance and proximity analysis are what let the
backend treat an accepting transcript as close to some valid encoded tensor.
Inside the unique-decoding radius this nearby codeword determines the tensor;
outside that radius one must reason with the stated list-decoding soundness
loss.  In either case, the received word itself is not promised to be error-free.
This is also the reason the Blaze compiler should see \(\BackendRMLE\) only as
an authentication boundary.  It submits tensors and coordinate handles, then
receives accept/reject answers for MLE claims.  It does not inspect the backend
codeword, choose backend proximity queries, or assume that verifier-facing
handles are literally error-free.

\begin{protocolbox}{Subroutine: \(\BackendRMLE\) Authentication}{fig:protocol-rmle-backend}
\scriptsize\raggedright
\noindent Public parameters: dimension \(M\); systematic foldable-code encoders
\(\Enc_i:\F^{\B^i}\to\F^{\B^{n_i}}\) for \(i\in\{0,\ldots,M\}\), with
\(\Enc_i(A)[0^{n_i-i}\Vert a]=A[a]\); the BaseFold folding rules from
Figure~\ref{fig:protocol-basefold-pcs-eval}; query count \(Q_{\mathrm{BF}}\);
and the coordinate-opening API \(\mathcal F_{\mathrm{MsgOracle}}\).

\smallskip
\noindent\(\underline{\BackendRMLE.\mathsf{Commit}(\prover:A\in\F^{\B^M},\ \verifier:C_A\in\Handle)\to A_{\mathrm{auth}}\in\Handle}\)

\smallskip
\begin{enumerate}[leftmargin=0.24in,label=\textbf{R\arabic*.},itemsep=0.06em,topsep=0.06em,parsep=0pt]
  \item \(\prover\): compute \(\pi_M=\Enc_M(A)\in\F^{\B^{n_M}}\).
  \item \(\prover\): let \(\pi_{\mathrm{ext}}\) be \(\pi_M\) restricted to
        \(\B^{n_M}\setminus(0^{n_M-M}\Vert\B^M)\).
  \item \(\prover\): \(C_{\mathrm{ext}}\gets
        \mathcal F_{\mathrm{MsgOracle}}.\mathsf{Submit}(\pi_{\mathrm{ext}})\).
  \item \(\prover\to\verifier\): \(C_{\mathrm{ext}}\).
  \item \(\prover,\verifier\): define \(C_M\) over \(\B^{n_M}\): on
        \(0^{n_M-M}\Vert a\), open \(C_A\) at \(a\); on
        \(b\in\B^{n_M}\setminus(0^{n_M-M}\Vert\B^M)\), open
        \(C_{\mathrm{ext}}\) at \(b\).
  \item \(\BackendRMLE\): return
        \(A_{\mathrm{auth}}=(C_M,C_{\mathrm{ext}},C_A,M)\), where \(C_M\) is the
        backend's top coordinate handle.
\end{enumerate}

\smallskip
\noindent\(\underline{\BackendRMLE.\mathsf{Eval}(\verifier:(A_{\mathrm{auth}},r\in\F^M,v\in\F),\ \prover:A\in\F^{\B^M})\to(\verifier:b\in\{0,1\})}\)

\smallskip
\begin{enumerate}[leftmargin=0.24in,label=\textbf{S\arabic*.},itemsep=0.06em,topsep=0.06em,parsep=0pt]
  \item \(\prover,\verifier\): run the backend RMLE evaluation proof for the
        claim \(\MLE(A)(r)=v\), using the commitment state in
        \(A_{\mathrm{auth}}\).
  \item \(\BackendRMLE\): internally, this proof may be instantiated by the
        BaseFold PCS evaluation protocol from
        Figure~\ref{fig:protocol-basefold-pcs-eval}, or by another IOPP/RMLE
        backend with the same authentication guarantee.
  \item \(\BackendRMLE\): return \((\verifier:b)\), where \(b\) is
        the result of that RMLE evaluation proof.
\end{enumerate}
\end{protocolbox}

\begin{protocolbox}{Subroutine: MLE Evaluation Batching}{fig:protocol-mle-eval-batch}
\scriptsize\raggedright
\noindent Public parameters: dimension \(M\), equality basis \(\chi\), and
degree-two \(M\)-variate sumcheck \(\mathsf{SC}_M\).

\smallskip
\noindent\(\underline{\mathsf{MLEBatch}(\verifier:(x_i\in\F^M,v_i\in\F)_{i\in[t]},\ \prover:A\in\F^{\B^M})\to(a\in\{0,1\},r_{\mathrm{batch}}\in\F^M,v_{\mathrm{batch}}\in\F)}\)

\smallskip
\begin{enumerate}[leftmargin=0.24in,label=\textbf{M\arabic*.},itemsep=0.06em,topsep=0.06em,parsep=0pt]
  \item \(\verifier\): sample \(\gamma\leftarrow\F\), then send
        \(\gamma\) to \(\prover\).
  \item \(\verifier\): for \(i\in[t]\), set \(\gamma_i\gets\gamma^{i-1}\).
  \item \(\verifier\): define
        \(W_\gamma(X)\gets\sum_{i\in[t]}\gamma_i\chi_X(x_i)\).
  \item \(\verifier\): set
        \(S_\gamma\gets\sum_{i\in[t]}\gamma_i v_i\).
  \item \(\prover\): set
        \(G_\gamma(X)\gets\MLE(A)(X)W_\gamma(X)\).
  \item \(\prover,\verifier\):
        \((a_{\mathrm{sc}},r_{\mathrm{batch}},T_{\mathrm{batch}})
        \gets\mathsf{SC}_M(\verifier:S_\gamma,\prover:G_\gamma)\).
  \item \(\verifier\): if \(a_{\mathrm{sc}}=0\), return
        \((0,0^M,0)\).
  \item \(\prover\to\verifier\): \(v_{\mathrm{batch}}\in\F\).
  \item \(\verifier\): if
        \(T_{\mathrm{batch}}\ne v_{\mathrm{batch}}W_\gamma(r_{\mathrm{batch}})\),
        return \((0,r_{\mathrm{batch}},v_{\mathrm{batch}})\).
  \item \(\verifier\): return \((1,r_{\mathrm{batch}},v_{\mathrm{batch}})\).
\end{enumerate}
\end{protocolbox}

\begin{protocolbox}{Implementation: \(\MLEAuth\) From Batching And BaseFold/RMLE}{fig:protocol-mle-auth-basefold}
\scriptsize\raggedright
\noindent Public parameters: field \(\F\); a maximum backend dimension \(D\);
an active lane set \(\mathcal L_0\subseteq\B^s\); a coordinate-opening API
\(\mathcal F_{\mathrm{MsgOracle}}\); the MLE batching subroutine from
Figure~\ref{fig:protocol-mle-eval-batch}; and the BaseFold-like RMLE subroutine
\(\BackendRMLE\) from Figure~\ref{fig:protocol-rmle-backend}, instantiated
with \(M=s+D\).

\smallskip
\noindent For \(h\in\F^{\B^d}\) with \(d\le D\), define the zero-padded tensor
\(\Pad_{d\to D}(h)\in\F^{\B^D}\) by
\[
  \Pad_{d\to D}(h)[a\Vert t]
  =
  \begin{cases}
    h[a], & t=0^{D-d},\\
    0, & t\ne0^{D-d},
  \end{cases}
  \qquad
  a\in\B^d,\quad t\in\B^{D-d}.
\]
Then
\[
  \MLE(\Pad_{d\to D}(h))(r\Vert0^{D-d})=\MLE(h)(r).
\]

\smallskip
\noindent\(\underline{\MLEAuth.\mathsf{Commit}(\prover:(h_q\in\F^{\B^{d_q}})_{q\in\mathcal L_0},\ \verifier:(d_q,C_q)_{q\in\mathcal L_0})\to A_{\mathrm{auth}}\in\Handle}\)

\smallskip
\begin{enumerate}[leftmargin=0.24in,label=\textbf{B\arabic*.},itemsep=0.06em,topsep=0.06em,parsep=0pt]
  \item \(\prover\): for each \(q\in\mathcal L_0\), compute
        \(\bar h_q\gets\Pad_{d_q\to D}(h_q)\in\F^{\B^D}\).
  \item \(\prover\): define the packed tensor
        \(A\in\F^{\B^s\times\B^D}\) by \(A[q,b]=\bar h_q[b]\) for
        \(q\in\mathcal L_0\), and \(A[q,b]=0\) for padded lanes
        \(q\in\B^s\setminus\mathcal L_0\).
  \item \(\prover,\verifier\): construct a packed coordinate handle \(C_A\).
        For \(q\in\mathcal L_0\) and \(b=a\Vert t\in\B^D\) with
        \(a\in\B^{d_q}\), opening \(C_A\) at \((q,b)\) routes to
        \(\mathcal F_{\mathrm{MsgOracle}}.\mathsf{Open}(C_q,a)\) when \(t=0^{D-d_q}\),
        and returns \(0\) otherwise.  Padded lanes also return \(0\).
  \item \(\prover,\verifier\): run the commit phase of
        \(\BackendRMLE\) on \((A,C_A)\), receiving backend handle
        \(A_{\mathrm{auth}}\).
  \item \(\MLEAuth\): return \(A_{\mathrm{auth}}\), with \(C_A\) and
        the public table \((q,d_q,C_q)_{q\in\mathcal L_0}\) attached to it.
\end{enumerate}

\smallskip
\noindent\(\underline{\MLEAuth.\mathsf{Eval}(\verifier:(A_{\mathrm{auth}},\mathcal E),\ \prover:(h_q)_{q\in\mathcal L_0})\to(\verifier:b\in\{0,1\})}\)

\smallskip
\noindent Here
\[
  \mathcal E=\{(q_i,r_i,v_i)\}_{i=1}^t,
  \qquad
  q_i\in\mathcal L_0,\quad r_i\in\F^{d_{q_i}},\quad v_i\in\F.
\]

\smallskip
\begin{enumerate}[leftmargin=0.24in,label=\textbf{O\arabic*.},itemsep=0.06em,topsep=0.06em,parsep=0pt]
  \item \(\prover\): use the packed tensor \(A\) determined by the commit
        phase from \((h_q)_{q\in\mathcal L_0}\).
  \item \(\verifier\): for \(i\in[t]\), set
        \(x_i\gets q_i\Vert r_i\Vert0^{D-d_{q_i}}\in\F^{s+D}\).
  \item \(\prover,\verifier\):
        \((a_{\mathrm{batch}},r_{\mathrm{batch}},v_{\mathrm{batch}})
        \gets\mathsf{MLEBatch}
        (\verifier:(x_i,v_i)_{i\in[t]},\prover:A)\).
  \item \(\verifier\): if \(a_{\mathrm{batch}}=0\), return
        \((\verifier:0)\).
  \item \(\prover,\verifier\): run the evaluation proof phase of
        \(\BackendRMLE\) on
        \((A_{\mathrm{auth}},r_{\mathrm{batch}},v_{\mathrm{batch}})\).
  \item \(\MLEAuth\): set \(b\gets1\) iff \(\BackendRMLE\) accepts;
        otherwise set \(b\gets0\).
  \item \(\MLEAuth\): return \((\verifier:b)\).
\end{enumerate}
\end{protocolbox}

This is our implementation model for the Blaze compiler boundary.  Blaze
supplies opaque MLIOP oracle tensors and later supplies evaluation requests.
The implementation pads and packs those oracles so that one BaseFold/RMLE
backend instance can bind the whole family.  Batching happens before the RMLE
\(\mathsf{Eval}\) proof:
handles are fixed first, then verifier randomness forms the batched claim.  A
batched virtual handle may still fan out to several underlying coordinate
openings when the backend queries it.  No RAA relation is injected into
BaseFold; the backend authenticates the packed tensor, while the RAA algebra is
carried entirely by the request list and the scalar checks that produced it.
